\chapter{LITERATURE REVIEW}
\label{ch:literature}
The literature connects two parts of the thesis question: why a voice resembling the listener might change an interaction, and why guided search is a useful setting in which to examine that possibility. Research on voice identity and self-similarity first motivates the comparison with a participant-preferred alternative. Work on search and auditory guidance then explains how spoken assistance can have a sustained role in the task. The review closes with the sensing and contextual requirements needed to make that comparison interpretable.

\section{Voice as a Social and Personalized Interface}
Research on machine-generated speech has considered personality, source attribution, and social presence~\cite{nass1993voices,nass2001personality,lee2005socialpresence}. These perspectives motivate treating voice as more than a neutral transport for words. They do not establish that every listener will form the same impression, or that a social impression necessarily improves task performance.

For this project, personalization has two meanings. Catalog personalization allows the participant to choose among existing voices. Self-similar personalization uses recorded speech to create a synthetic voice resembling that participant. Both forms may produce familiarity or liking, but they operate through different routes. Preference is selected directly in the catalog condition, whereas perceived resemblance must be checked in the self-similar condition. The manipulation is therefore not adequately described as personalized versus generic.

The bibliography also includes work on voice-assistant personality customization and similarity~\cite{snyder2023onevoice}. Personality similarity is conceptually different from acoustic self-resemblance. A listener can prefer a speaking style without hearing it as their own voice, and a close acoustic resemblance can be experienced as uncomfortable. The planned manipulation check asks about perceived voice similarity, while the experience battery separately asks about liking, naturalness, familiarity, and uncanniness.

A related AR navigation system by Yang et al. combines Building Information Modeling with retrieval-augmented generation and compares directional arrows alone with an embodied guide accompanied by the same arrows~\cite{yang2025embodiednavigation}. Its comparison concerns the presentation of navigation assistance, not self-similar voice. It provides a useful contrast for the present framing: the question here is how the guiding voice is experienced within a shared assistance mechanism, rather than whether a new reasoning architecture or visible agent improves navigation. The walking task, destination queries, and embodied interface in that system differ from the seated hidden-target task used here.

\section{Self-Similar Voice and Embodied Collaboration}
Self-similar voice provides the empirical motivation for this study, while auditory guidance supplies the setting in which that voice is encountered. Two particularly relevant studies examine resemblance in different roles: an avatar speaking as the player and a virtual collaborator speaking to the participant. The present study extends this question to a guiding voice without a visible agent body.

Kao et al. analyzed 657 participants in an online programming game, comparing self-similar and assigned dissimilar avatar voices while also manipulating pitch~\cite{kao2021selfsimilar}. Self-similar voice was associated with more puzzles completed, longer voluntary play, and higher similarity identification, competence, relatedness, and immersion. Direct differences on the reported intrinsic-motivation subscales were not significant. These findings motivate examining several dimensions of voice experience rather than assuming that resemblance improves enjoyment or every measure of motivation.

The meaning of performance also changes across tasks. Longer voluntary play in Kao et al.'s game can indicate continued engagement; a longer fixed search trial here can indicate difficulty. Follow My Voice will therefore interpret completion time alongside workload, motivation, and reported reliance, without treating speed as a complete account of the value of self-similar voice. This is a transfer of the research question, not a prediction that the direction of every outcome will replicate.

Guo et al. used a 25-participant, within-subject design crossing appearance and voice self-similarity during virtual jigsaw collaboration~\cite{guo2024doppelganger}. Self-similar voices increased likability and believability ratings while also increasing eeriness. Voice did not significantly affect co-presence or perceived intelligence, although appearance affected those measures. These results support keeping favorable and adverse reactions separate and distinguishing voice effects from effects of a visible character. They do not establish a search-speed advantage.

Both studies compared self-similar voices with assigned alternatives; Follow My Voice instead compares the clone with a participant-preferred catalog voice~\cite{kao2021selfsimilar,guo2024doppelganger}. The substantive question is consequently whether resemblance to oneself offers a different experience from a voice one has already chosen. Recognition of the clone, liking of the voice, and willingness to follow its instructions need not coincide. Keeping these constructs separate makes a mixed outcome informative rather than contradictory.

The distinction between embodiment of the participant and embodiment of the agent is important. Participants wear a headset and use a tracked controller, so the interaction involves physical sensing and action. The guiding agent, however, is not represented by a co-present animated character. Measures adapted from an embodied collaborator must therefore be described as adaptations. Their use does not establish that the same factor structure or interpretation holds for a disembodied voice.

Self-voice applications in daily reminders and emotional self-voice interventions extend the range of contexts represented in the bibliography~\cite{kim2024selfvoice,fang2025emotional}. They motivate asking whether participants can imagine uses for their own synthetic voice. They do not establish that navigation guidance will be more effective or that participants will experience it as inner speech. The debrief consequently begins with open questions before introducing probes about eeriness, self-talk, and possible applications.

\section{Naturalness, Familiarity, and Uncanniness}
Voice resemblance is not the only acoustic difference that listeners may notice. Synthesis can change articulation, prosody, loudness, duration, or perceived naturalness. Research on vocal uncanniness provides a reason to keep these qualities visible rather than subsume them under similarity~\cite{diel2024uncanny}. A clone that is recognizable but difficult to understand is not equivalent to a recognizable and intelligible clone.

The experiment will use the same semantic library and provider for both conditions. This is an implementation control rather than proof of equal acoustics, because prepared clips can differ in duration and pronunciation even when their text matches. Pilot checks will therefore inspect intelligibility and comfortable playback levels across catalog and cloned voices, and the actual provider, model, voice, and library identifiers will identify the configuration each participant heard.

Custom experience items and the appealing, eerie, and familiar items selected from Guo et al. will remain separate, since combining them would imply a measurement model not established for this study. A participant may report a voice as both familiar and eerie, and preserving the individual item directions keeps that response interpretable.

\section{Workload, Motivation, and Agent Perception}
The possibility of both welcome and uncomfortable self-resemblance makes outcome selection part of the research argument. Task speed describes efficiency, while workload, motivation, and agent perception address different aspects of receiving help. The following measures retain those distinctions so that the comparison can reveal a tradeoff rather than require one voice to rank higher on every dimension.

Zagermann et al. discuss eye tracking as a means of investigating cognitive load in visual computing~\cite{zagermann2016workload}. That motivates considering gaze-derived measures, but does not license interpreting an unvalidated ray-hit count as workload. The current confirmatory plan uses subjective workload dimensions and task performance. Pupil, fixation, and saccade analyses remain contingent on separate acquisition and event-validation work.

NASA-TLX supplies six workload dimensions~\cite{hart1988nasatlx}; the project's 0--10 response range and unweighted complete mean are documented adaptations. The Intrinsic Motivation Inventory distinguishes interest/enjoyment from related subscales such as effort, pressure, and value~\cite{ryan1982control,mcauley1989imi,csdtIMI}, and reporting the subscales separately preserves that distinction. The instrument appendix specifies the current wording and reverse keys so that scoring can be reproduced.

Agent perception measures are drawn through Guo et al., with attribution to their originating measures~\cite{guo2024doppelganger,biocca2001networked,moussawi2019perception}. Co-presence, perceived intelligence, and anthropomorphism concern participant attributions, so a high rating does not establish sentience or decision correctness. Conversely, deterministic validity checks do not establish that a participant found the agent natural. The design therefore needs both kinds of evidence while keeping their meanings separate.


\section{Visual Search and Attention Guidance in Extended Reality}
To study a guiding voice, the task must give its instructions a meaningful role. Object search provides that role by requiring the participant to connect guidance with a location and then an intended object. The search literature helps specify this setting; it does not supply the answer to the voice comparison.

Search in a surrounding display differs from search on a single front-facing screen because relevant objects may lie outside the visible field. Turning changes what can be inspected, while eye movements refine attention within the visible region. Work on eye and head movements across an extended field of view is therefore relevant to the task geometry~\cite{stein2024eyehead}. In the present study, this motivates recording both head motion and gaze rather than treating one as a complete substitute for the other.

Several approaches to directing attention in immersive environments inform this design. Grogorick et al. address subtle gaze guidance, while HiveFive concerns immersion-preserving attention guidance~\cite{grogorick2017subtle,lange2020hivefive}. Ariza and colleagues investigate vibrotactile assistance and proximity-based multimodal feedback for selection~\cite{ariza2017vibrotactile,ariza2018proximity}. These works establish relevant design alternatives without implying that all forms of guidance are equivalent: a cue that changes visual salience, a vibration that indicates proximity, and a spoken spatial instruction impose different perceptual and interpretive demands.

RippleVision provides a recent visual counterpart to gaze-contingent spoken guidance. Kudnick et al. describe brightness ripples restricted to a cone between a point of interest and the current gaze location, with opposite brightness modulation in the two eyes~\cite{kudnick2026ripplevision}. This technique description is supported by the published abstract; detailed study methods and comparative outcomes remain outside the present synthesis pending full-text review. Its relevance here is conceptual: assistance can adapt to where a participant looks without specifying the target through spoken identity. It also raises a design question for voice guidance: how much assistance is useful before it becomes intrusive? That question motivates considering participant experience alongside task performance; it does not establish that visual and spoken cues have equivalent effects. Follow My Voice holds its assistance policy common across voice conditions to investigate self-similar voice, rather than comparing guidance modalities.

Binetti et al. examine visual and auditory cues for locating out-of-view objects in head-mounted augmented reality~\cite{binetti2021outofview}, and Warden et al. focus on target cue type and location in augmented-reality search~\cite{warden2022visual}. Together, these studies show why target visibility and cue availability belong in the method description. Follow My Voice does not open with a color-and-shape description of the target; instead, the sequence of directions gradually narrows the participant's search.

This hidden-target design creates an important boundary. A participant who already knows the target can continue searching without further assistance. Here, the participant must obtain enough spatial information to infer which object is intended. The voice may consequently be heard more often and may be more consequential than in a task where it supplements an already known target. That is a reason to describe the task precisely, not a reason to assume stronger voice effects in advance.

\section{Gaze-Contingent Auditory Guidance}
Auditory guidance is not a single technique. A sound can encode gaze error through loudness, indicate a location through spatial rendering, or communicate an instruction through words. Only the last of these necessarily carries the verbal and vocal identity cues central to this thesis. The following studies inform the shared search design while leaving the self-similar voice comparison open.

Losing et al. tested gaze-contingent sonification in a small desktop study with ten participants: the loudness of a flute tone varied with gaze-to-target distance, reaching silence at the target~\cite{losing2014guiding}. Benefits depended on the search task, with the clearest improvement in the difficult number-search condition. This is evidence for linking auditory feedback to current gaze, rather than evidence for any particular spoken phrase. Our adaptation uses short spatial instructions instead of a continuously varying tone so that participants encounter the assigned voice throughout search. Its frequency and timing remain design choices to validate, not benefits established by the sonification study.

Kwok et al.'s Two-Step Gaze Guidance evaluated 69 participants using a three-monitor display and eye tracking, rather than a VR headset~\cite{kwok2022twostep}. Initial spatial-audio orientation followed by finer guidance shortened search times relative to the other tested conditions; accuracy did not differ significantly across conditions. Their implementation also buffered transitions between guidance stages. The transferable idea is a staged change in the information a cue supplies: first orient toward a region, then refine within it. Follow My Voice adapts that structure to wall entry and within-wall directions; the cited study does not validate this headset, spoken wording, full-surround geometry, or transition thresholds.

Binetti et al. provide a complementary headset-based example using visual arrows and spatialized alarm sounds to locate out-of-view targets~\cite{binetti2021outofview}. Adding sound mitigated response-time costs when visual cues and targets occupied different depth planes, but did not remove the corresponding accuracy cost. The study had no sound-only condition. Its relevance is therefore conditional support for an auditory channel when visual guidance is constrained, not proof that speech is superior to visual cues. It also motivates retaining selection errors separately from completion time in the present voice comparison.

Zhang et al.'s 96-participant AR book-search study more directly concerns verbal instructions~\cite{zhang2023seehear}. Its audio hint specified an approximate shelf once at task onset; it was not repeated gaze-contingent speech. In the sub-study without gaze-playback feedback, hint conditions outperformed the control but did not significantly differ from one another. Combined visual/audio advantages over either modality appeared in the second task of the feedback sub-study. These results support considering task completion and workload together, while making cue timing and feedback exposure explicit. They do not establish that continuous spoken guidance or self-similar delivery will be better.

\paragraph{Design synthesis for the voice comparison.}
The resulting design uses a common coarse-to-fine guidance policy to make the two voice conditions comparable. The agent first helps locate the relevant wall, then narrows the search within it; the participant confirms the inferred target with the right trigger. The target's initially hidden identity gives the voice an ongoing instructional role, unlike a single announcement supplementing a known target. Whether hearing that role in one's own synthetic voice changes performance, comfort, motivation, or reliance remains the empirical question.

Guidance during search, acknowledgment of a qualified gaze or controller contact, and feedback after selection are recorded separately. This separation also distinguishes the researcher's Trial Review tool from Zhang et al.'s participant-facing gaze-playback intervention~\cite{zhang2023seehear}. Reviewing a recording to audit delivery does not expose a participant to a learning intervention. Similarly, a success chime cannot explain an earlier change in direction within the same trial, although it may affect subsequent experience. No participant playback intervention is added to the current protocol.

The current study does not include an unguided or gaze-unaware control; both voice conditions receive the same gaze-contingent policy. Earlier project documents that considered a factorial combination of voice and gaze contingency are historical, and their interpretation cannot be transferred to this two-condition design. Any observed difference will therefore concern the assigned voice experience within the shared assistance system.

\section{Context and Short-Term Memory in XR Agents}
EmBARDiment uses gaze-informed contextual memory to support interaction with an embodied XR agent~\cite{bovo2025embardiment}. That approach motivates making recent interaction evidence available to the assistance policy, although Follow My Voice implements a narrower, bounded interaction history. This history records recent spatial progress, actual speech, and guidance transitions, and its retention limits are implementation choices rather than thresholds taken from the cited system; it does not provide long-term autobiographical memory or an inferred model of the participant's mental state. The guiding voice also has no visible agent body, so the present system does not inherit the embodiment of the cited interface.

Memory-Augmented AR Agents and GenLARP represent additional directions in the bibliography~\cite{choi2025memory,yu2025genlarp}. The former is a conceptual framework and the latter a short system description; neither supplies a validated timing threshold for this application. User-aware virtual agents and proactive AR assistance also offer context for designing state-dependent behavior~\cite{gunawardhana2024useraware,li2025satori}. These works position a design problem, rather than prove that this implementation's memory mechanism improves performance.

Short-term history is operationally useful when the same current geometry can have different meanings. A target to the left may be one the participant has not yet approached, or one they have just passed while turning right. A recently delivered cue can make another identical cue unnecessary. History distinguishes these cases only if it records reliable observations and actual playback. Holding this mechanism constant across voices means the study cannot estimate its independent effect.

Feedback specificity offers a further distinction between adaptive assistance and generic encouragement. Zimmerer et al. examined performance-specific feedback from a virtual exercise coach and found changes in selected agent-perception measures without a corresponding effect on motivation~\cite{zimmerer2025feedback}. For Follow My Voice, this distinction motivates describing whether a cue reports observed progress, corrects direction, acknowledges qualified gaze or pointing, or confirms a successful selection. These functions remain common across voice conditions, so the study will compare voice configurations rather than estimate an independent effect of feedback specificity. Results from embodied exercise coaching do not establish an effect on search time in the present task.

\section{Eye-Tracking Quality and Interpretation}
Eye-tracking quality involves more than a nominal sampling frequency. The reporting guidance of Dunn et al. and the data-quality discussion of Niehorster et al. motivate reporting acquisition, validity, spatial performance, and missingness~\cite{dunn2024reporting,niehorster2026dataquality}. The system must establish how many distinct source observations arrive, how old they are when used, and how error varies across the search space.

Validation work on the Vive Pro Eye and tools such as GazeMetrics are methodological resources~\cite{schuetz2022vive,sipatchin2021usability,adhanom2020gazemetrics}. They cannot substitute for validation on the HTC Vive Focus Vision because device model, fit, calibration, coordinate conversion, and task geometry can differ. The protocol therefore requires independent targets across heights, boundaries, distances, and head orientations, and these validation targets must be separate from targets used to estimate an offset correction.

Drewes et al. compare gaze during locomotion in virtual reality and the real world~\cite{drewes2021locomotion}. This study is seated and uses head and chair rotation rather than free walking, so locomotion findings cannot be imported as direct estimates of its gaze behavior. Head motion will be logged because it participates in search, while claims about eye events remain bounded by the available measurements.

Salvucci and Goldberg discuss identifying fixations and saccades from eye-tracking protocols~\cite{salvucci2000fixations}. An interaction region around an object serves a different purpose. The 100-ms vicinity qualification in this system is an interface acknowledgment rule, not a validated fixation classifier. The thesis will report gaze contacts and qualified vicinity events under their own names.

\section{Comparison of the Closest Research Strands}
Table~\ref{tab:related-work-map} compares the selected work by what it contributes to this design and what it leaves unresolved. It is a focused synthesis of the preceding review, not a systematic-review inventory or a ranking of systems. The literature does not supply one interchangeable baseline: spatial guidance, voice resemblance, context management, and measurement quality answer different questions.

\begin{table}[htbp]\centering
\caption[Related research and boundaries of transfer]{Selected research strands and their role in Follow My Voice. These are methodological connections, not evidence that the present implementation reproduces the cited effects.}
\label{tab:related-work-map}
\small\begin{tabularx}{\linewidth}{@{}p{0.26\linewidth}XX@{}}\toprule
Research strand & Relevance to this project & Boundary of transfer \\\midrule
Self-similar voice~\cite{kao2021selfsimilar,guo2024doppelganger} & Motivates the voice manipulation and both favorable and adverse experience measures. & Educational or embodied collaboration tasks do not establish a search-speed benefit here. \\
Auditory and staged gaze guidance~\cite{losing2014guiding,kwok2022twostep} & Supports investigating gaze-dependent, coarse-to-fine assistance. & Does not validate this headset, phrase policy, or gate durations. \\
AR search and task feedback~\cite{zhang2023seehear} & Separates in-task cues from post-task feedback. & Known-target task and feedback manipulation differ from this hidden-target voice comparison. \\
Context-aware XR agents~\cite{bovo2025embardiment} & Motivates explicit recent context for state-dependent assistance. & Does not establish an independent benefit of this bounded memory. \\
Eye-tracking reporting and validation~\cite{dunn2024reporting,niehorster2026dataquality,adhanom2020gazemetrics} & Informs source timing, validity, spatial validation, and gap reporting. & Reporting guidance and other-device results are not Focus Vision accuracy measurements. \\\bottomrule
\end{tabularx}
\end{table}

The resulting contribution is a particular combination of task and comparison. The agent knows an initially unnamed target, the participant receives adaptive spatial speech, and explicit right-trigger selection remains separate from gaze. The empirical manipulation changes the voice configuration while retaining that assistance mechanism. This combination defines the scope of RQ1--RQ3; it does not imply that voice familiarity, memory, and gaze contingency have been independently manipulated.

\section{Synthesis and Research Gap}
The selected literature motivates a focused question about self-similar voice: how does guidance spoken in a synthetic version of the participant's own voice compare with guidance spoken in a preferred catalog voice? Prior voice studies make both favorable and uncomfortable reactions plausible, while audio-guidance studies inform a shared task through which those reactions and performance can be examined. Gaze contingency, staged refinement, and decision support serve that comparison; they are not additional experimental treatments.

The selected work leaves the present comparison unresolved: what changes when the voice guiding an initially unnamed, surrounding object-search task resembles its listener rather than being a preferred catalog choice? The design combines a consequential use for speech with a comparator that already accommodates preference. Its contribution will be to establish how those two voice configurations differ in this setting, including outcomes that favor neither configuration consistently. The next chapter describes the shared apparatus that makes the comparison possible.
